% ======================================================================
\section{Conclusions}\label{sec:conclusions}
% ======================================================================
% SLOT (partially drafted). Inputs: all sections; the paragraphs on the collinear, running-coupling
% and finite terms must be written after sections dglap_rc and finite (P1-P6) are filled in.

We have computed the single inclusive gluon spectrum at mid rapidity in the collision of a dilute projectile with a dense target to next-to-leading order in the Color Glass Condensate, in pure Yang--Mills theory. The calculation uses the evolution operator $\Omega$ of the projectile, written as a normal-ordered series in the soft gluon operators with coefficients that are functionals of the non-commuting valence charge density, and reduces the spectrum to products of these coefficients and of Wilson lines. We have organised the $O(g^4)$ cross section into three groups of terms and fifteen rows.

The rapidity-enhanced terms have the structure required by high-energy factorisation. The logarithm of the longitudinal cut-offs splits at the momentum of the produced gluon. The coefficient of $\ln(\vee/k^+)$ is the BFKL evolution of the projectile charge correlator inserted into the LO spectrum, and the coefficient of $\ln(k^+/\Lambda)$ is the JIMWLK Hamiltonian acting on the target operator of the LO spectrum, for an arbitrary ordered projectile correlator. Terms with four charges cancel, as required by the coherent-state structure of the LO wave function, and the terms with three charges cancel in the symmetrised sense; with non-commuting charges, their reordering produces part of the two-charge structures of both evolutions. The operator ordering of the charges, in particular of the three-charge term of the $O(g^3)$ wave function, is essential for these cancellations.

\pending{To be completed: collinear term and fragmentation, running coupling, finite terms, open points, outlook.}
